% English translation of chapter2.tex lines 721--910.
% Source: Methods of Algebra, Volume 2, commit
% 9a5803ff77dd3257484cb177f851a73770a59dd3 (CC BY 4.0).
% stable-unit-id: o014.aljabr2.chapter2.direct-sum-decomposition

% segment-id: o014.li2.chapter2.direct-sum-decomposition.h001
\section{Direct-Sum Decompositions}\label{sec:direct-sum}
\hypertarget{o014.aljabr2.chapter2.direct-sum-decomposition}{ }

% segment-id: o014.li2.chapter2.direct-sum-decomposition.p001
We begin by discussing how morphisms between direct sums are represented by
matrices. The argument is exactly the same as in the module-theoretic case
\cite[\S 6.3]{Li1}; only a brief review is given here.

% segment-id: o014.li2.chapter2.direct-sum-decomposition.p002
Fix an additive category $\mathcal{A}$. Consider objects
$X_1,\ldots,X_n$ and $X'_1,\ldots,X'_m$ of $\mathcal{A}$. The direct sum
$\bigoplus_{i=1}^nX_i$ introduced in
Definition~\ref{def:additive-category} comes with a family of morphisms
$X_j\xrightarrow{\iota_j}\bigoplus_{i=1}^nX_i\xrightarrow{p_j}X_j$.
Likewise, $\bigoplus_{i=1}^mX'_i$ has morphisms $p'_j,\iota'_j$, and so on.
The universal properties of products and coproducts give
% segment-id: o014.li2.chapter2.direct-sum-decomposition.q001
\begin{align*}
	\Hom \left( \bigoplus_{j=1}^n X_j, \bigoplus_{i=1}^m X'_i \right) & \xrightarrow[\sim]{\phi \mapsto (p'_i \phi)_i } \prod_{i=1}^m \Hom\left(\bigoplus_{j=1}^n X_j , X'_i \right) \\
	& \xrightarrow[\sim]{(p'_i \phi)_i \mapsto (p'_i \phi \iota_j)_{i, j}} \prod_{j=1}^n \prod_{i=1}^m  \Hom(X_j, X'_i).
\end{align*}
Thus every morphism
$\phi:\bigoplus_{j=1}^nX_j\to\bigoplus_{i=1}^mX'_i$ corresponds to the
matrix
% segment-id: o014.li2.chapter2.direct-sum-decomposition.q002
\[ \mathcal{M}(\phi) = (\phi_{ij})_{\substack{1 \leq i \leq m \\ 1 \leq j \leq n}} = \begin{pmatrix}
	\phi_{11} & \cdots & \phi_{1n} \\
	\vdots & \ddots & \vdots \\
	\phi_{m1} & \cdots & \phi_{mn}
\end{pmatrix}, \quad \phi_{ij} := p'_i \phi \iota_j: X_j \to X'_i. \]

% segment-id: o014.li2.chapter2.direct-sum-decomposition.p003
Conversely, every matrix
$\mathcal{M}=(\phi_{ij})_{\substack{1\leq i\leq m\\1\leq j\leq n}}$
formed from morphisms $\phi_{ij}:X_j\to X'_i$ uniquely determines a
morphism $\phi$ such that $\mathcal{M}=\mathcal{M}(\phi)$. Composition of
morphisms agrees with matrix multiplication:
% segment-id: o014.li2.chapter2.direct-sum-decomposition.q003
\[ \mathcal{M}(\psi\phi) = \mathcal{M}(\psi) \mathcal{M}(\phi) \]
where multiplication of matrix entries is given by composition
$\Hom(X'_j,X''_i)\times\Hom(X_k,X'_j)\to\Hom(X_k,X''_i)$.

% segment-id: o014.li2.chapter2.direct-sum-decomposition.le001
\begin{lemma}\label{prop:isomorphism-matrix}
	Given a family of morphisms $\phi_i:X_i\to X'_i$ in $\mathcal{A}$,
	$i=1,\ldots,n$, set
	$\phi:=(\phi_1,\ldots,\phi_n):\bigoplus_{i=1}^nX_i\to
	\bigoplus_{i=1}^nX'_i$. Then $\phi$ is an isomorphism if and only if each
	$\phi_i$ is an isomorphism.
\end{lemma}
% segment-id: o014.li2.chapter2.direct-sum-decomposition.pf001
\begin{proof}
	The implication $\impliedby$ is clear. For $\implies$, write
	write the matrix $\mathcal{M}(\phi^{-1})$ as $(\psi_{ij})_{i,j}$. Then
	\[ \mathcal{M}(\phi^{-1}) \mathcal{M}(\phi) = \begin{pmatrix}
		\psi_{11} \phi_1 & \cdots & \psi_{1n} \phi_n \\
		\vdots & \ddots & \vdots \\
		\psi_{n1} \phi_1 & \cdots & \psi_{nn} \phi_n
	\end{pmatrix} = \begin{pmatrix}
		\identity_{X_1} & & \\
		& \ddots & \\
		& & \identity_{X_n}
	\end{pmatrix}\]
	so $\phi_1,\ldots,\phi_n$ all have left inverses. The same argument shows
	that they all have right inverses.
\end{proof}

% segment-id: o014.li2.chapter2.direct-sum-decomposition.p004
Continue to consider a direct-sum decomposition
$X\simeq\bigoplus_{i=1}^nX_i$, where $n\geq1$ and $X_i\neq0$ for every
$i$. It gives morphisms $\iota_i:X_i\to X$ and $p_i:X\to X_i$. For each
$1\leq i\leq n$, set $e_i:=\iota_ip_i\in\End(X)$. These morphisms have the
following properties. Recall that an \emph{idempotent} in the ring $\End(X)$
is an element $e\in\End(X)$ satisfying $e^2=e$; when no confusion can arise,
we also call $e$ an idempotent in $\mathcal{A}$.
\index{idempotent}

% segment-id: o014.li2.chapter2.direct-sum-decomposition.l001
\begin{enumerate}[(E1)]
	% segment-id: o014.li2.chapter2.direct-sum-decomposition.i001
	\item Every $e_i$ is idempotent:
	$e_i^2=(\iota_ip_i)(\iota_ip_i)=\iota_i(p_i\iota_i)p_i=\iota_ip_i$.
	% segment-id: o014.li2.chapter2.direct-sum-decomposition.i002
	\item Orthogonality: $i\neq j\implies
	e_ie_j=\iota_ip_i\iota_jp_j=0$.
	% segment-id: o014.li2.chapter2.direct-sum-decomposition.i003
	\item The equality $\sum_{i=1}^ne_i=1$ holds in the ring $\End(X)$.
\end{enumerate}

% segment-id: o014.li2.chapter2.direct-sum-decomposition.p005
From now on we regard the direct summands $X_i$ as subobjects of $X$ and
write the decomposition as an equality $X=\bigoplus_{i=1}^nX_i$.
Conversely, if $\mathcal{A}$ is abelian, a family of idempotents
$e_1,\ldots,e_n\in\End(X)$ satisfying (E1)--(E3) determines subobjects of
$X$ by
% segment-id: o014.li2.chapter2.direct-sum-decomposition.q005
\[ X_i := \Ker\left( \sum_{j \neq i} e_j \right) = \Image(e_i), \quad i=1, \ldots, n. \]
On the one hand they come with monomorphisms $X_i\xrightarrow{\iota_i}X$;
on the other, $e_i:X\to X$ factors uniquely as
$X\xrightarrow{p_i}X_i\xrightarrow{\iota_i}X$. This is precisely the data
required for a direct sum.

% segment-id: o014.li2.chapter2.direct-sum-decomposition.p006
Passing from the idempotents $e_1,\ldots,e_n$ to the direct summands
$X_1,\ldots,X_n$ does not require all the properties of an abelian category.
It is enough to assume that $\mathcal{A}$ is an $\cate{Ab}$-category with a
zero object in which every idempotent has a kernel. Such a category is called
a \emph{Karoubian category} or a \emph{pseudo-abelian category}. The exercises
of this chapter discuss this further.
\index{category!Karoubian}

% segment-id: o014.li2.chapter2.direct-sum-decomposition.pr001
\begin{proposition}\label{prop:idempotent-decomposition}
	Let $X\neq0$ be an object of an abelian category (or, more generally, a
	Karoubian category) $\mathcal{A}$, and let $n\in\Z_{\geq1}$. The
	constructions above give a bijection
	\[\begin{tikzcd}[row sep=small]
		\left\{\begin{gathered}
			\text{direct-sum decompositions}\\[-2pt]
			X = \bigoplus_{i=1}^n X_i
		\end{gathered}\right\}
		\arrow[leftrightarrow, r, "1:1"] &
		\left\{\begin{gathered}
			(e_i)_{i=1}^n \in \End(X)^n:\\[-2pt]
			\text{satisfying (E1)--(E3)}
		\end{gathered}\right\}.
	\end{tikzcd}\]
	Reordering the direct summands $X_1,\ldots,X_n$ corresponds to reordering
	the idempotents $e_1,\ldots,e_n$.
\end{proposition}
% segment-id: o014.li2.chapter2.direct-sum-decomposition.pf002
\begin{proof}
	See \cite[Proposition 6.12.4]{Li1}.
\end{proof}

% segment-id: o014.li2.chapter2.direct-sum-decomposition.p007
In the special case $n=2$, direct-sum decompositions are related to a special
class of short exact sequences, called split short exact sequences.

% segment-id: o014.li2.chapter2.direct-sum-decomposition.pr002
\begin{proposition}\label{prop:split-ses}
	\index{short exact sequence!split}
	For a short exact sequence
	$0\to X'\xrightarrow{f}X\xrightarrow{g}X''\to0$ in an abelian category,
	the following statements are equivalent.
	\begin{enumerate}[(i)]
		% segment-id: o014.li2.chapter2.direct-sum-decomposition.i004
		\item There is an $s:X''\to X$ such that $gs=\identity_{X''}$.
		% segment-id: o014.li2.chapter2.direct-sum-decomposition.i005
		\item There is an $r:X\to X'$ such that $rf=\identity_{X'}$.
		% segment-id: o014.li2.chapter2.direct-sum-decomposition.i006
		\item There is a diagram
		$\begin{tikzcd}
			X' \arrow[yshift=-0.7ex, r, "f"'] & X \arrow[yshift=0.7ex, l, "r"'] \arrow[yshift=-0.7ex, r, "g"'] & X'' \arrow[yshift=0.7ex, l, "s"']
		\end{tikzcd}$
		for which $X\simeq X'\oplus X''$; see the review of biproducts in
		\S\ref{sec:Ker-Coker}.
		% segment-id: o014.li2.chapter2.direct-sum-decomposition.i007
		\item The map $g_*:\Hom(S,X)\to\Hom(S,X'')$, sending $\varphi$ to
		$g\varphi$, is surjective for every object $S$.
		% segment-id: o014.li2.chapter2.direct-sum-decomposition.i008
		\item The map $f^*:\Hom(X,S)\to\Hom(X',S)$, sending $\psi$ to
		$\psi f$, is surjective for every object $S$.
	\end{enumerate}
	If any of these conditions holds, the short exact sequence
	$0\to X'\to X\to X''\to0$ is called \emph{split}. The morphism
	$s:X''\to X$ in (i), the morphism $r:X\to X'$ in (ii),% 
	\footnote{Editorial correction (O014-C024): the source writes
	$r:X'\to X$ here, but (ii), the equality $rf=\identity_{X'}$, and the
	subsequent construction all require $r:X\to X'$. The type-correct direction
	is used here.}
	and the isomorphism $X\rightiso X'\oplus X''$ in (iii), denoted by $\Phi$,
	correspond as follows:
	\begin{compactitem}
		% segment-id: o014.li2.chapter2.direct-sum-decomposition.i009
		\item from $\Phi$ to $s$, take the composite
		$X''\xrightarrow{\iota_2}X'\oplus X''\xrightarrow{\Phi^{-1}}X$;
		% segment-id: o014.li2.chapter2.direct-sum-decomposition.i010
		\item from $s$ to $r$, factor $\identity_X-sg:X\to X$ uniquely as
		$X\xrightarrow{r}X'\xrightarrow{f}X$; this is the desired $r$;
		% segment-id: o014.li2.chapter2.direct-sum-decomposition.i011
		\item from $r$ to $\Phi$, take $(r,g):X\to X'\oplus X''$, which is an
		isomorphism.
	\end{compactitem}

	Moreover, for a given $s$ or its corresponding $r$, the orthogonal
	idempotents corresponding to the direct-sum decomposition in (iii) are
	\[ e' := fr = \identity_X - sg, \quad e'' := sg = \identity_X - fr. \]
\end{proposition}
% segment-id: o014.li2.chapter2.direct-sum-decomposition.pf003
\begin{proof}
	For modules, see \cite[Proposition 6.8.5]{Li1}. Its proof uses only
	algebraic operations in $\Hom$ sets and the characterization of biproducts,
	all of which also hold in an abelian category, so the argument carries over
	without change. We omit the details.
\end{proof}

% segment-id: o014.li2.chapter2.direct-sum-decomposition.p008
In the setting of Proposition~\ref{prop:split-ses}, given $f:X'\to X$ and
$g:X\to X''$, an $s$ satisfying $gs=\identity_{X''}$ is called a
\emph{section} of $g$, while an $r$ satisfying $rf=\identity_{X'}$ is called
a \emph{retraction} of $f$; these terms come from topology. The existence of
a section or retraction implies, respectively, that $g$ is an epimorphism or
$f$ is a monomorphism, but the converses do not hold.
\index{section}
\index{retraction}

% segment-id: o014.li2.chapter2.direct-sum-decomposition.p009
Notice also that the five statements (i)--(v) of the proposition are
collectively self-dual.

% segment-id: o014.li2.chapter2.direct-sum-decomposition.d001
\begin{definition}\label{def:indecomposable}\index{object!indecomposable}
	Let $S$ be a nonzero object of an abelian category $\mathcal{A}$. If
	$S=S'\oplus S''$ implies $S'=0$ or $S''=0$, then $S$ is called
	\emph{indecomposable}.
\end{definition}

% segment-id: o014.li2.chapter2.direct-sum-decomposition.co001
\begin{corollary}\label{prop:indecomposable-criterion}
	A nonzero object $X$ of an abelian category is indecomposable if and only
	if
	\[ \forall e \in \End(X), \quad e^2 = e \iff (e = 0 \;\vee\; e = 1). \]
\end{corollary}
% segment-id: o014.li2.chapter2.direct-sum-decomposition.pf004
\begin{proof}
	Apply Proposition~\ref{prop:idempotent-decomposition}.
\end{proof}

% segment-id: o014.li2.chapter2.direct-sum-decomposition.p010
We wish to study decompositions $X=X_1\oplus\cdots\oplus X_n$, where
$X\neq0$ and every $X_i$ is indecomposable. There are two questions:
existence and uniqueness. For modules this is the content of the
Krull--Remak--Schmidt theorem; see \cite[Corollary 6.12.9]{Li1}. For a
general abelian category, some preparation is needed. We follow the approach
of \cite{Kr15}.

% segment-id: o014.li2.chapter2.direct-sum-decomposition.d002
\begin{definition}\label{def:bichain}
	\index{bi-chain condition}
	Let $\mathcal{A}$ be any category. A \emph{bi-chain} in $\mathcal{A}$ is
	data $(X_n,\alpha_n,\beta_n)_{n=0}^\infty$, where the $X_n$ are objects
	and
	$\begin{tikzcd}
		X_n \arrow[twoheadrightarrow, r, shift left, "\alpha_n"] & X_{n+1} \arrow[hookrightarrow, l, shift left, "\beta_n"]
	\end{tikzcd}$
	are morphisms between them, with $\alpha_n$ an epimorphism and $\beta_n$ a
	monomorphism for $n\in\Z_{\geq0}$. An object $X$ of $\mathcal{A}$ is said
	to satisfy the \emph{bi-chain condition} if, for every bi-chain
	$(X_n,\alpha_n,\beta_n)_{n=0}^\infty$ with $X_0=X$, both $\alpha_n$ and
	$\beta_n$ are isomorphisms for $n\gg0$.
\end{definition}

% segment-id: o014.li2.chapter2.direct-sum-decomposition.p011
The following result generalizes \cite[Lemma 6.11.5]{Li1}.
% segment-id: o014.li2.chapter2.direct-sum-decomposition.le002
\begin{lemma}[Fitting's lemma in an abelian category]\label{prop:Abel-cat-Fitting}
	Let $\mathcal{A}$ be an abelian category, let $X$ be a nonzero object
	satisfying the bi-chain condition, and let $f\in\End(X)$.
	\begin{enumerate}[(i)]
		% segment-id: o014.li2.chapter2.direct-sum-decomposition.i012
		\item For $n\gg0$, $\Ker(f^n)$ and $\Image(f^n)$ are independent of
		$n$; denote them by $\Ker(f^\infty)$ and $\Image(f^\infty)$,
		respectively. Then
		$X=\Ker(f^\infty)\oplus\Image(f^\infty)$.
		% segment-id: o014.li2.chapter2.direct-sum-decomposition.i013
		\item If $X$ is indecomposable, then $f$ is either invertible or
		nilpotent in the ring $\End(X)$.
	\end{enumerate}
\end{lemma}
% segment-id: o014.li2.chapter2.direct-sum-decomposition.pf005
\begin{proof}
	For $n\in\Z_{\geq0}$, define $X_n:=\Image(f^n)$, so $X_0=X$.
	Proposition~\ref{prop:Images-f-gf} and $\Coim\rightiso\Image$ give an
	epimorphism $\alpha_n:X_n\to X_{n+1}$ induced by $f$ and a monomorphism
	$\beta_n:X_{n+1}\to X_n$. Thus we obtain a bi-chain
	$(X_n,\alpha_n,\beta_n)_{n=0}^\infty$.

	For $n\gg0$, $\alpha_n$ and $\beta_n$ are isomorphisms. Hence the
	subobject $\Image(f^\infty):=\Image(f^n)$ of $X$, for $n\gg0$, is well
	defined; write its monomorphism as
	$\iota:\Image(f^\infty)\hookrightarrow X$. Likewise,
	$\Ker(f^n)=\Ker[X\to\Image(f^n)]$ is a subobject $\Ker(f^\infty)$
	independent of $n$ for $n\gg0$.

	For $n\gg0$, the morphism
	$\alpha_{2n-1}\cdots\alpha_n:X_n\to X_{2n}$ is invertible. Denote its
	inverse by
	$\psi:\Image(f^\infty)\rightiso\Image(f^\infty)$ and set
	\[ p := \psi f^n: X \to \Image(f^\infty), \quad \Ker(p) = \Ker(f^\infty) \qquad (n \gg 0). \]
	Since $p\iota=\identity_{\Image(f^\infty)}$, the endomorphism
	$e:=\iota p\in\End(X)$ is idempotent, with
	$\Image(e)=\Image(f^\infty)$ and $\Ker(e)=\Ker(f^\infty)$.
	Proposition~\ref{prop:idempotent-decomposition} gives the decomposition in
	(i).

	If $X$ is indecomposable, then either $\Image(f^\infty)=0$ and
	$\Ker(f^\infty)=X$, in which case $f$ is nilpotent, or
	$\Image(f^\infty)=X$ and $\Ker(f^\infty)=0$, in which case $f$ is
	invertible. This proves (ii).
\end{proof}

% segment-id: o014.li2.chapter2.direct-sum-decomposition.p012
Recall from \cite[Definition 6.12.3]{Li1} that if $S$ is a ring and
$S\smallsetminus S^\times$ is a two-sided ideal, then $S$ is called
\emph{local}.
\index{local ring}

% segment-id: o014.li2.chapter2.direct-sum-decomposition.co002
\begin{corollary}\label{prop:indecomposable-local-ring}
	Let $X$ be a nonzero object in an abelian category satisfying the bi-chain
	condition. Then $X$ is indecomposable if and only if $\End(X)$ is a local
	ring.
\end{corollary}
% segment-id: o014.li2.chapter2.direct-sum-decomposition.pf006
\begin{proof}
	Lemma~\ref{prop:Abel-cat-Fitting} shows that if $X$ is indecomposable, every
	element of $\End(X)$ is either nilpotent or invertible. The remainder of
	the argument concerns only the ring structure of $\End(X)$ and is the same
	as \cite[Lemma 6.12.6]{Li1}.
\end{proof}

% segment-id: o014.li2.chapter2.direct-sum-decomposition.p013
We can now state the version of the Krull--Remak--Schmidt theorem for abelian
categories.
% segment-id: o014.li2.chapter2.direct-sum-decomposition.th001
\begin{theorem}[M.\ Atiyah]\label{prop:Krull-Schmidt-gen}
	\index{Krull--Remak--Schmidt theorem}
	Let $X$ be a nonzero object of an abelian category.
	\begin{enumerate}[(i)]
		% segment-id: o014.li2.chapter2.direct-sum-decomposition.i014
		\item If $X$ satisfies the bi-chain condition, then there exist
		$n\in\Z_{\geq1}$ and indecomposable subobjects $X_1,\ldots,X_n$ such
		that $X=\bigoplus_{i=1}^nX_i$, and every $\End(X_i)$ is a local ring.
		% segment-id: o014.li2.chapter2.direct-sum-decomposition.i015
		\item Suppose that $X$ has decompositions $\bigoplus_{i=1}^nX_i$ and
		$\bigoplus_{j=1}^mX'_j$, where every $X_i$ and $X'_j$ is
		indecomposable and every $\End(X_i)$ and $\End(X'_j)$ is local. Then
		$n=m$, and there is a bijection $\sigma$ of $\{1,\ldots,n\}$ such that
		$X_i\simeq X'_{\sigma(i)}$ for every $i$.
	\end{enumerate}
\end{theorem}
% segment-id: o014.li2.chapter2.direct-sum-decomposition.pf007
\begin{proof}
	The main point is to prove that an $X$ satisfying the bi-chain condition
	has a decomposition as in (i). Everything else involves only algebraic
	operations in the rings $\End$ and is the same as
	\cite[Theorem 6.12.8]{Li1}. We therefore assume the bi-chain condition in
	what follows.

	If $X=Y\oplus Z$ with $Y\neq0$, then $Y$ also satisfies the bi-chain
	condition. Indeed, given a bi-chain
	$(Y_n,\alpha_n,\beta_n)_{n=0}^\infty$ with $Y_0=Y$, take
	$X_n:=Y_n\oplus Z$, $\tilde{\alpha}_n:=\alpha_n\oplus\identity_Z$, and
	$\tilde{\beta}_n:=\beta_n\oplus\identity_Z$. This gives a bi-chain with
	$X_0=X$, while Lemma~\ref{prop:isomorphism-matrix} shows that
	$\tilde{\alpha}_n$ (or $\tilde{\beta}_n$) is an isomorphism if and only if
	$\alpha_n$ (or $\beta_n$) is.

	If $X$ is already indecomposable,
	Corollary~\ref{prop:indecomposable-local-ring} shows that $\End(X)$ is
	local, proving (i). If (i) were false for $X$, there would be nonzero
	subobjects $X_1,Y_1$ such that $X=X_1\oplus Y_1$ and (i) remained false
	for $X_1$. Repeating the process on $X_1$ gives
	$X_1=X_2\oplus Y_2$, and iteration
	produces a bi-chain $(X_n,\alpha_n,\beta_n)_{n=0}^\infty$, where
	$\alpha_n:X_n\twoheadrightarrow X_{n+1}$ and
	$\beta_n:X_{n+1}\hookrightarrow X_n$ arise from
	$X_n=X_{n+1}\oplus Y_{n+1}$. None is an isomorphism, and $X_0:=X$,
	contradicting the bi-chain condition. This proves the assertion.
\end{proof}

% segment-id: o014.li2.chapter2.direct-sum-decomposition.p014
In applications, the key issue is knowing when the bi-chain condition holds.
Here is one sufficient condition.

% segment-id: o014.li2.chapter2.direct-sum-decomposition.pr003
\begin{proposition}
	Let $X$ be an object of finite length in an abelian category
	(Definition~\ref{def:finite-length-object}). Then $X$ satisfies the
	bi-chain condition.
\end{proposition}
% segment-id: o014.li2.chapter2.direct-sum-decomposition.pf008
\begin{proof}
	Let $(X_n,\alpha_n,\beta_n)_{n=0}^\infty$ be a bi-chain with $X_0=X$.
	Then
	\[
		\left(\Image(\beta_0\cdots\beta_n)\right)_{n=0}^\infty
	\]
	is a descending chain in the partially ordered set $\mathrm{Sub}_X$. Since all
	$\beta_n$ are monomorphisms, for $n\gg0$ the Artinian condition gives a
	commutative diagram
	\[\begin{tikzcd}
		X_{n+1} \arrow[d, "\beta_n"'] \arrow[r, "\sim"] & \Image(\beta_0 \cdots \beta_n) \arrow[d, "\sim" sloped] \arrow[hookrightarrow, r] & X_0 \\
		X_n \arrow[r, "\sim"] & \Image(\beta_0 \cdots \beta_{n-1}) \arrow[hookrightarrow, ru] &
	\end{tikzcd}\]
	and $\beta_n$ is an isomorphism. Similarly, apply the Noetherian condition
	to the ascending chain
	$\left(\Ker(\alpha_n\cdots\alpha_0)\right)_{n=0}^\infty$ and use the fact
	that every $\alpha_n$ is an epimorphism. For $n\gg0$ this gives a
	commutative diagram with exact rows
	\[\begin{tikzcd}
		0 \arrow[r] & \Ker(\alpha_{n-1} \cdots \alpha_0) \arrow[d, "\sim" sloped] \arrow[r] & X_0 \arrow[r, "{\alpha_{n-1} \cdots \alpha_0}" inner sep=0.7em] \arrow[equal, d] & X_n \arrow[d, "{\alpha_n}"] \arrow[r] & 0 \\
		0 \arrow[r] & \Ker(\alpha_n \cdots \alpha_0) \arrow[r] & X_0 \arrow[r, "{\alpha_n \cdots \alpha_0}"' inner sep=0.7em] & X_{n+1} \arrow[r] & 0
	\end{tikzcd}\]
	and $\alpha_n$ is an isomorphism.
\end{proof}

% segment-id: o014.li2.chapter2.direct-sum-decomposition.p015
The exercises introduce further ways to deduce the bi-chain condition.
